\documentclass[11pt,a4paper]{article}
\usepackage[utf8]{inputenc}
\usepackage[T1]{fontenc}
\usepackage{lmodern}
\usepackage[french]{babel}
\usepackage[margin=2.5cm]{geometry}
\usepackage{xcolor}
\usepackage{graphicx}
\usepackage{booktabs,longtable,tabularx,array}
\usepackage{enumitem}
\usepackage{fancyhdr}
\usepackage{listings}
\usepackage{url}
\usepackage{tikz}
\usetikzlibrary{arrows.meta,positioning,shapes.geometric,fit}
\usepackage[colorlinks=true,linkcolor=blue!50!black,urlcolor=blue!50!black]{hyperref}

% ---------- À MODIFIER ----------
\newcommand{\titreprojet}{Agent DevOps}
\newcommand{\sousTitre}{Surveillance et fiabilisation de la chaîne de livraison logicielle}
\newcommand{\membres}{Hazeel \\ \textit{[Nom du membre 2]} \\ \textit{[Nom du membre 3]}}
\newcommand{\encadrant}{\textit{[Nom de l'enseignant]}}
\newcommand{\annee}{2026--2027}
\newcommand{\version}{1.0}
\newcommand{\datedoc}{23 novembre 2026}
% --------------------------------

\definecolor{accent}{HTML}{1F4E79}
\setlist{itemsep=2pt,topsep=4pt}
\setlength{\parindent}{0pt}\setlength{\parskip}{6pt}
\renewcommand{\arraystretch}{1.25}
\lstset{basicstyle=\ttfamily\footnotesize,frame=single,breaklines=true,
        backgroundcolor=\color{gray!8},xleftmargin=4pt,xrightmargin=4pt,
        literate={é}{{\'e}}1 {è}{{\`e}}1 {à}{{\`a}}1 {ê}{{\^e}}1 {ç}{{\c c}}1}
\newcommand{\code}[1]{\texttt{#1}}

\pagestyle{fancy}\fancyhf{}
\fancyhead[L]{\small\titreprojet\ -- Cahier des charges}
\fancyhead[R]{\small Version \version}
\fancyfoot[C]{\thepage}
\renewcommand{\headrulewidth}{0.4pt}

\begin{document}

% ================= PAGE DE GARDE =================
\begin{titlepage}
\centering
{\large École Marocaine des Sciences de l'Ingénieur (EMSI)}\\[0.2cm]
{\normalsize Filière Informatique et Réseaux -- Parcours IA et Data -- 4\textsuperscript{e} année}\\[0.2cm]
{\normalsize Module : IA et Développement}\\[2.5cm]
{\color{accent}\rule{\linewidth}{1.2pt}}\\[0.5cm]
{\Huge\bfseries\color{accent} CAHIER DES CHARGES}\\[0.4cm]
{\LARGE \titreprojet}\\[0.3cm]
{\large\itshape \sousTitre}\\[0.3cm]
{\color{accent}\rule{\linewidth}{1.2pt}}\\[2cm]
\begin{tabular}{rl}
\textbf{Réalisé par :} & \begin{tabular}[t]{l}\membres\end{tabular} \\[0.8cm]
\textbf{Encadrant :} & \encadrant \\[0.4cm]
\textbf{Année universitaire :} & \annee \\[0.4cm]
\textbf{Version / date :} & \version\ -- \datedoc
\end{tabular}
\vfill
{\small Catégorie : Exploitation et infrastructure}
\end{titlepage}

\tableofcontents
\newpage

% ================= 1 =================
\section{Introduction et contexte}
Dans les équipes de développement modernes, chaque modification du code traverse une chaîne
d'intégration et de livraison continues (CI/CD) : installation des dépendances, tests, contrôle
de qualité, construction, déploiement. Lorsqu'une étape échoue, un ingénieur doit lire des
journaux (\emph{logs}) souvent longs et bruités, retrouver l'étape fautive, comprendre la cause,
puis corriger ou relancer. Cette tâche répétitive consomme du temps d'exploitation et retarde les livraisons.

Le projet \titreprojet{} vise à automatiser la première ligne de ce travail : un agent qui
\textbf{observe} les pipelines, \textbf{diagnostique} les échecs, \textbf{propose} un correctif
et, dans un cadre de sécurité strict, \textbf{agit} sur la chaîne de livraison.

Ce document définit le besoin, le périmètre, les exigences, l'architecture, les choix
technologiques, la stratégie de validation et le planning du projet. Il est réalisé dans le
cadre du module « IA et Développement » et sert de référence pour les jalons GitHub (semaines 1 à 3).

% ================= 2 =================
\section{Problématique et objectifs}
\subsection{Problématique}
\emph{Comment réduire le temps de diagnostic et de reprise après un échec CI/CD, sans
compromettre la sécurité de l'infrastructure par des actions automatiques non maîtrisées~?}

\subsection{Objectif général}
Surveiller et fiabiliser la chaîne de livraison logicielle grâce à un agent capable d'analyser les
logs, de détecter les erreurs CI/CD, de proposer une correction et de lancer certaines actions
automatisées sous contrôle.

\subsection{Objectifs spécifiques et indicateurs de réussite}
\begin{longtable}{>{\bfseries}p{1.2cm}p{7.3cm}p{6cm}}
\toprule
\textbf{ID} & \textbf{Objectif} & \textbf{Indicateur mesurable} \\
\midrule\endhead
O1 & Identifier correctement la cause d'un échec (catégorie et étape) & Exactitude de catégorie $\geq 85\,\%$ sur le jeu d'évaluation \\
O2 & Recommander l'action adaptée (relance, PR, rollback, aucune) & Exactitude d'action $\geq 85\,\%$ \\
O3 & Ne jamais exécuter d'action sensible sans validation humaine & 0 rollback exécuté sans approbation (testé) \\
O4 & Garantir la traçabilité complète & 100\,\% des décisions et actions présentes dans le journal d'audit \\
O5 & S'abstenir en cas de doute plutôt que de se tromper & Un log inconnu produit une abstention et une escalade \\
O6 & Livrer un dépôt reproductible et documenté & Installation et lancement réussis sur un clone neuf \\
\bottomrule
\end{longtable}

% ================= 3 =================
\section{Périmètre}
\subsection{Inclus}
\begin{itemize}
\item Lecture des logs de pipelines CI/CD (GitHub Actions en priorité ; mode simulé pour la démonstration).
\item Classification des échecs : dépendance manquante, test cassé, configuration erronée, erreur
d'authentification, échec réseau transitoire, erreur de lint, erreur de build, échec de déploiement, mémoire insuffisante.
\item Proposition d'un correctif en langage naturel, avec les preuves extraites du log.
\item Actions encadrées : relance d'un job, rollback d'un déploiement, ouverture d'une pull request corrective.
\item Validation humaine pour les actions sensibles, journal d'audit, API REST et interface web.
\end{itemize}
\subsection{Exclus}
\begin{itemize}
\item Fusion automatique de code dans la branche principale.
\item Gestion ou modification des secrets (jetons, clés) : l'agent escalade toujours vers un humain.
\item Prise en charge de tous les fournisseurs CI/CD (extension envisagée en perspectives).
\item Correction automatique de bugs fonctionnels complexes : l'agent propose, l'humain corrige.
\end{itemize}

% ================= 4 =================
\section{Acteurs et besoins}
\begin{longtable}{p{3.2cm}p{5.6cm}p{5.7cm}}
\toprule
\textbf{Acteur} & \textbf{Rôle} & \textbf{Besoin principal} \\
\midrule\endhead
Développeur & Pousse du code, subit les échecs de pipeline & Comprendre vite pourquoi le build échoue et comment corriger \\
Ingénieur DevOps / exploitation & Responsable des déploiements & Réduire le temps de reprise, valider ou refuser les actions sensibles \\
Responsable sécurité & Garant des règles d'accès & Garantir que l'agent reste dans ses droits et que tout est tracé \\
Agent (système) & Analyse et agit selon les droits accordés & Décider dans le cadre de la politique de sécurité \\
\bottomrule
\end{longtable}

% ================= 5 =================
\section{Exigences fonctionnelles}
Priorités selon la méthode MoSCoW : \textbf{M} (indispensable), \textbf{S} (important), \textbf{C} (souhaitable).
\begin{longtable}{>{\bfseries}p{1.2cm}p{9.6cm}p{1.2cm}p{1.5cm}}
\toprule
\textbf{ID} & \textbf{Exigence} & \textbf{Prio.} & \textbf{Jalon} \\
\midrule\endhead
EF1 & Récupérer les logs d'un run CI/CD (par identifiant ou par texte collé) & M & S2 \\
EF2 & Masquer les secrets présents dans les logs avant stockage, affichage ou envoi à un LLM & M & S2 \\
EF3 & Classer l'échec par catégorie et par étape, avec un score de confiance & M & S2 \\
EF4 & Extraire les lignes de preuve justifiant le diagnostic & M & S2 \\
EF5 & Proposer un correctif adapté à la catégorie & M & S2 \\
EF6 & Recommander une action : relance, PR, rollback ou aucune & M & S2 \\
EF7 & Appliquer une politique de sécurité (risque, seuil de confiance, nombre maximal de relances) & M & S2 \\
EF8 & Exiger une validation humaine pour toute action sensible (rollback) ou peu fiable & M & S2 \\
EF9 & Exécuter les actions autorisées (relance, rollback, PR), avec mode simulation (\code{DRY\_RUN}) & M & S3 \\
EF10 & Journaliser chaque décision et chaque action (qui, quand, quoi, pourquoi) & M & S2 \\
EF11 & S'abstenir et escalader quand aucun signal fiable n'est trouvé & M & S3 \\
EF12 & Exposer une API REST documentée (santé, analyse, incidents, approbation, audit) & M & S2 \\
EF13 & Fournir une interface web : analyser, consulter les incidents, approuver ou rejeter & S & S2 \\
EF14 & Générer une explication en langage naturel via un LLM (facultatif, mode dégradé sans LLM) & S & S3 \\
EF15 & Ouvrir réellement une pull request corrective sur GitHub & C & S3 \\
EF16 & Surveiller les déploiements en continu (interrogation périodique) & C & S3 \\
\bottomrule
\end{longtable}

% ================= 6 =================
\section{Exigences non fonctionnelles}
\begin{longtable}{>{\bfseries}p{3.4cm}p{11.1cm}}
\toprule
\textbf{Catégorie} & \textbf{Exigence} \\
\midrule\endhead
Sécurité & Principe du moindre privilège ; aucun secret dans le dépôt (\code{.env} ignoré) ; secrets masqués dans les logs ; l'agent n'accède jamais aux secrets ; protection contre le \emph{path traversal} sur les identifiants de run. \\
Traçabilité & Journal d'audit persistant et horodaté de chaque diagnostic, demande d'approbation, approbation, rejet et action. \\
Fiabilité & Une panne du LLM ou du fournisseur CI/CD ne doit pas faire tomber l'agent : mode dégradé et erreurs explicites. \\
Performance & Diagnostic par règles en moins de 1 seconde pour un log de 20\,000 caractères ; logs tronqués en conservant la fin. \\
Maintenabilité & Code typé et testé, contrôle automatique (Ruff, pytest) à chaque push, règles de détection ajoutables en une entrée. \\
Reproductibilité & Versions figées dans \path{requirements.txt}, installation en trois commandes, mode simulé sans compte externe. \\
Utilisabilité & Interface simple : diagnostic, correctif, risque et boutons d'approbation visibles sur une seule page. \\
Erreurs & Format d'erreur stable de l'API : \code{\{"error": \{"code", "message"\}\}}. \\
\bottomrule
\end{longtable}

% ================= 7 =================
\section{Cas d'usage principaux}
\subsection{CU1 -- Échec dû à une dépendance manquante}
\textbf{Acteur} : agent, développeur. \textbf{Déclencheur} : un job de tests échoue.
\begin{enumerate}
\item L'agent récupère les logs du run, masque les secrets et conserve la fin du log.
\item Il détecte la signature « module introuvable », catégorie \code{missing\_dependency}, confiance élevée.
\item Il propose : « ajouter la dépendance X à \path{requirements.txt} » et recommande l'ouverture d'une PR.
\item La politique autorise l'action (risque moyen, confiance suffisante) ; la PR est ouverte (ou simulée) et tracée.
\end{enumerate}
\textbf{Alternative} : confiance sous le seuil, l'action passe en attente d'approbation.

\subsection{CU2 -- Déploiement défaillant en production}
\begin{enumerate}
\item Le déploiement échoue (healthcheck, \emph{CrashLoopBackOff}).
\item L'agent diagnostique \code{deploy\_failure} et recommande un rollback.
\item Le rollback étant sensible, l'incident passe en \code{awaiting\_approval}.
\item L'ingénieur DevOps approuve ou rejette dans l'interface ; la décision est tracée avec son auteur.
\end{enumerate}

\subsection{CU3 -- Échec transitoire (réseau)}
L'agent reconnaît un délai d'attente dépassé, recommande une relance ; elle est exécutée
automatiquement dans la limite de \code{MAX\_RETRIES}, puis une validation humaine est requise.

\subsection{CU4 -- Log inconnu ou erreur d'authentification}
Aucune signature fiable : l'agent \textbf{s'abstient}, explique pourquoi et escalade. Pour une erreur d'authentification,
il ne tente aucune action et invite à renouveler le secret côté plateforme.

% ================= 8 =================
\section{Architecture technique}
\subsection{Vue d'ensemble}
L'architecture est organisée en couches : une API REST (FastAPI) expose l'agent ; l'agent
orchestre des outils (récupération des logs, analyse, explication, politique, actions) ; un
connecteur CI/CD abstrait le fournisseur (simulé ou GitHub) ; un journal d'audit persiste les événements.

\begin{center}
\begin{tikzpicture}[node distance=0.9cm and 1.1cm,
 box/.style={draw=accent,fill=accent!8,rounded corners,minimum width=2.6cm,minimum height=0.9cm,align=center,font=\small},
 ext/.style={draw=gray,fill=gray!10,rounded corners,minimum width=2.6cm,minimum height=0.9cm,align=center,font=\small},
 arr/.style={-{Stealth[length=2mm]},thick,color=accent}]
\node[ext] (ui) {Interface\\React};
\node[box,right=of ui] (api) {API FastAPI};
\node[box,right=of api] (agent) {Agent\\(boucle)};
\node[ext,right=of agent] (ci) {CI/CD\\GitHub Actions};
\node[box,below=of agent] (an) {Analyseur\\règles + abstention};
\node[box,left=of an] (pol) {Politique\\de sécurité};
\node[ext,right=of an] (llm) {LLM\\(optionnel)};
\node[box,below=of pol] (act) {Actions\\relance/rollback/PR};
\node[box,below=of an] (aud) {Journal d'audit\\SQLite};
\draw[arr] (ui) -- (api); \draw[arr] (api) -- (agent); \draw[arr] (agent) -- (ci);
\draw[arr] (agent) -- (an); \draw[arr] (agent) -- (pol); \draw[arr] (agent) -- (llm);
\draw[arr] (pol) -- (act); \draw[arr] (agent) to[bend left=18] (aud);
\end{tikzpicture}
\end{center}

\subsection{Boucle de l'agent}
\begin{enumerate}
\item \code{fetch\_logs} : récupération du log (run identifié ou texte fourni).
\item \code{analyze\_logs} : masquage des secrets, troncature, détection par règles, score de confiance ; abstention si aucun signal.
\item \code{explain\_with\_llm} (optionnel) : explication en langage naturel ; en cas d'indisponibilité, le diagnostic par règles est conservé.
\item \code{policy} : calcul du risque et décision (automatique, validation humaine ou aucune action).
\item Exécution de l'action ou attente d'approbation ; chaque étape est ajoutée à la trace et au journal d'audit.
\end{enumerate}

\subsection{Stratégie d'analyse en deux niveaux}
Le premier niveau (règles) est rapide, gratuit, déterministe et testable ; il couvre les cas
fréquents. Le second niveau (LLM) enrichit l'explication mais ne décide jamais seul d'une
action : la décision d'agir reste du ressort de la politique de sécurité. Cette séparation
limite les risques d'hallucination et de comportement imprévisible.

% ================= 9 =================
\section{Choix technologiques}
\begin{longtable}{p{2.6cm}p{3.6cm}p{8.3cm}}
\toprule
\textbf{Domaine} & \textbf{Choix} & \textbf{Justification} \\
\midrule\endhead
Langage backend & Python 3.12 & Écosystème IA, lisibilité, tests rapides à écrire. \\
API & FastAPI + Pydantic & Validation automatique, documentation interactive \code{/docs}, typage fort. \\
Serveur & Uvicorn & Serveur ASGI standard pour FastAPI. \\
Client HTTP & httpx & Appels GitHub et LLM, délais d'attente configurables, testable. \\
Analyse des logs & Expressions régulières (règles) & Déterministe, explicable, sans entraînement ; extensible. \\
IA générative & API d'un LLM (Claude) via HTTP & Explication et suggestion en langage naturel ; facultatif et remplaçable. \\
CI/CD & GitHub Actions & Dépôt déjà sur GitHub, API de logs, relance et \emph{workflow\_dispatch} pour le rollback. \\
Persistance & SQLite (bibliothèque standard) & Aucune installation, suffisant pour le journal d'audit. \\
Frontend & React + Vite & Démarrage rapide, composants simples, proxy de développement vers l'API. \\
Tests et qualité & pytest, Ruff, GitHub Actions & Tests automatiques et contrôle de style à chaque push. \\
Conteneurs (option) & Docker & Packaging et exécution reproductible si le temps le permet. \\
\bottomrule
\end{longtable}

\textbf{Alternatives écartées.} LangChain/LangGraph (dépendance lourde pour une boucle simple),
entraînement d'un modèle de classification (jeu de données insuffisant en une semaine),
PostgreSQL (surdimensionné), Jenkins (installation lourde).

% ================= 10 =================
\section{Politique de sécurité des actions}
\begin{longtable}{p{3cm}p{1.8cm}p{9.7cm}}
\toprule
\textbf{Action} & \textbf{Risque} & \textbf{Règle appliquée} \\
\midrule\endhead
\code{rerun\_job} & Faible & Automatique si confiance $\geq$ seuil et relances $<$ \code{MAX\_RETRIES} ; sinon validation humaine. \\
\code{open\_pr} & Moyen & Automatique si confiance $\geq$ seuil ; la PR est relue par un humain, jamais fusionnée par l'agent. \\
\code{rollback} & Élevé & Validation humaine \textbf{toujours} obligatoire. \\
Secrets, authentification & -- & Aucune action : escalade vers un humain. \\
Confiance sous le seuil & -- & Toute action passe en validation humaine. \\
\bottomrule
\end{longtable}
Mesures complémentaires : mode simulation \code{DRY\_RUN} activé par défaut ; option \code{AUTO\_ACTIONS}
pour désactiver toute action automatique ; jeton GitHub aux droits minimaux ; masquage des secrets ;
journal d'audit consultable ; identifiants de run assainis.

% ================= 11 =================
\section{Modèle de données et API}
\subsection{Données principales}
\begin{longtable}{p{3cm}p{11.5cm}}
\toprule
\textbf{Objet} & \textbf{Champs principaux} \\
\midrule\endhead
Diagnostic & catégorie, étape, confiance, explication, correctif proposé, action recommandée, preuves, abstention \\
Décision & risque, action automatique autorisée, validation requise, raison \\
Incident & identifiant, date, run, environnement, diagnostic, décision, statut, explication LLM, trace des outils, résultat d'action \\
Événement d'audit & horodatage, incident, acteur, événement, détails \\
\bottomrule
\end{longtable}
Statuts d'un incident : \code{resolved\_auto}, \code{awaiting\_approval}, \code{executed}, \code{rejected}, \code{no\_action}.

\subsection{Routes de l'API}
\begin{longtable}{p{4.7cm}p{9.8cm}}
\toprule
\textbf{Route} & \textbf{Rôle} \\
\midrule\endhead
\code{GET /health} & État du service et mode courant \\
\code{GET /runs} & Liste des runs disponibles \\
\code{POST /analyze} & Analyse d'un run (\code{run\_id}) ou d'un log (\code{log\_text}) \\
\code{GET /incidents}, \code{GET /incidents/\{id\}} & Consultation des incidents \\
\code{POST /incidents/\{id\}/approve} et \code{/reject} & Validation humaine \\
\code{GET /audit} & Journal d'audit \\
\bottomrule
\end{longtable}

% ================= 12 =================
\section{Stratégie de tests et d'évaluation}
\subsection{Tests automatisés}
Tests unitaires de l'analyseur (catégories, abstention, log vide ou très long, masquage des secrets),
de la politique de sécurité (seuils, rollback, limite de relances) et de l'API (statuts, format
d'erreur, validation humaine, journal d'audit, \emph{path traversal}). Les tests s'exécutent à chaque push via l'intégration continue.

\subsection{Jeu d'évaluation}
Un jeu de 15 à 20 cas (log, catégorie attendue, action attendue) est versionné dans \path{data/eval_cases.json}.
Le script \path{scripts/evaluate.py} calcule l'exactitude de catégorie et d'action et écrit le résultat dans \path{docs/evaluation.md}.
Les cas initiaux seront complétés par des logs réels issus des dépôts du groupe afin de limiter le biais d'un jeu écrit par les concepteurs.

\subsection{Critères de réussite}
\begin{itemize}
\item Exactitude de catégorie et d'action $\geq 85\,\%$ sur le jeu d'évaluation.
\item Aucune action sensible exécutée sans validation (vérifié par test).
\item Journal d'audit complet pour 100\,\% des incidents.
\item Cas limites couverts : log vide ou trop long, fournisseur indisponible, confiance faible (abstention).
\end{itemize}

% ================= 13 =================
\section{Organisation et planning}
\subsection{Workflow Git}
Une branche par jalon (par exemple \code{s2-3-api}), pull request relue par un autre membre, fusion
dans \code{main} uniquement si les tests passent ; messages de commit de la forme
\code{S2-3 : ajouter la route principale et ses tests} ; un tag par étape (\code{s1}, \code{s2}, \code{v1.0.0}) ;
chaque membre est auteur principal d'au moins un jalon par étape, depuis son compte personnel.

\subsection{Planning sur une semaine}
Le planning ci-dessous condense les trois semaines de référence en sept jours ; les 18 jalons
et les trois tags sont conservés. \emph{Ce rythme doit être confirmé avec l'enseignant.}
\begin{longtable}{p{2.7cm}p{8cm}p{3.8cm}}
\toprule
\textbf{Jour} & \textbf{Travail} & \textbf{Jalons} \\
\midrule\endhead
Lun. 23/11 & Cahier des charges, dépôt, environnement, squelette, données, architecture, fiche 1 & S1-1 à S1-6, tag \code{s1} \\
Mar. 24/11 & Cœur de l'agent (boucle, deux outils) et jeu d'évaluation & S2-1, S2-2 \\
Mer. 25/11 & API (santé, analyse, erreurs stables, tests) & S2-3 \\
Jeu. 26/11 & Interface reliée à l'API, intégration continue (tests, Ruff) & S2-4, S2-5 \\
Ven. 27/11 & Fiche 2, README à jour, tag \code{s2} & S2-6 \\
Sam. 28/11 & Cas limites, évaluation finale, installation sur clone neuf & S3-1 à S3-3 \\
Dim. 29/11 & Documentation, compte rendu, support de soutenance, release & S3-4 à S3-6, tag \code{v1.0.0} \\
\bottomrule
\end{longtable}

\subsection{Répartition indicative}
\begin{longtable}{p{3.5cm}p{11cm}}
\toprule
\textbf{Rôle} & \textbf{Responsabilités} \\
\midrule\endhead
Backend et agent & Analyseur, politique de sécurité, actions, connecteur CI/CD, audit \\
API et qualité & Routes, validation, tests, intégration continue, évaluation \\
Frontend et documentation & Interface React, README, schémas, compte rendu et soutenance \\
\bottomrule
\end{longtable}
Les rôles seront ajustés selon la taille du groupe ; chacun reste auteur d'au moins un jalon par étape.

% ================= 14 =================
\section{Risques et mitigations}
\begin{longtable}{p{4.2cm}p{1.8cm}p{8.5cm}}
\toprule
\textbf{Risque} & \textbf{Gravité} & \textbf{Mitigation} \\
\midrule\endhead
Délai d'une semaine trop court & Élevée & Squelette fourni, périmètre « M » d'abord, fonctionnalités « C » en dernier \\
Action automatique dangereuse & Élevée & Politique stricte, \code{DRY\_RUN} par défaut, validation humaine, tests dédiés \\
Fuite d'un secret dans le dépôt ou le log & Élevée & \code{.gitignore}, \code{.env.example}, masquage, remplacement immédiat de toute clé poussée \\
Hallucination du LLM & Moyenne & LLM limité à l'explication ; décision par règles et politique \\
Indisponibilité de GitHub ou du LLM & Moyenne & Mode simulé, mode dégradé, erreurs explicites \\
Jeu d'évaluation biaisé & Moyenne & Ajouter de vrais logs du groupe, documenter les limites \\
Conflits Git dans l'équipe & Faible & Une branche par jalon, revue croisée, \code{main} toujours exécutable \\
\bottomrule
\end{longtable}

% ================= 15 =================
\section{Livrables}
\begin{itemize}
\item Dépôt GitHub complet : code, tests, intégration continue, historique montrant la contribution de chacun.
\item README : objectif, installation, lancement, exemples, limites.
\item Documentation : architecture, évaluation, \code{CHANGELOG.md}, fiches de suivi des trois semaines.
\item Compte rendu de projet et support de soutenance.
\item Version finale étiquetée \code{v1.0.0} avec notes de version.
\end{itemize}

% ================= 16 =================
\section{Perspectives}
Prise en charge d'autres fournisseurs (GitLab CI, Jenkins), ouverture réelle de PR avec correctif
généré et testé, apprentissage à partir des incidents passés, notification (Slack, e-mail),
tableau de bord de fiabilité (taux d'échec, temps moyen de reprise).

\appendix
\section{Structure du dépôt}
\begin{lstlisting}
agent-devops/
|-- README.md  requirements.txt  .env.example  .gitignore  CHANGELOG.md
|-- .github/workflows/ci.yml
|-- src/        main.py agent.py analyzer.py policy.py actions.py
|               cicd.py llm.py audit.py schemas.py config.py
|-- tests/      test_analyzer.py test_policy.py test_api.py
|-- data/       logs/ (runs d'exemple)  eval_cases.json
|-- scripts/    evaluate.py
|-- docs/       architecture.md evaluation.md equipe.md suivi/ cahier-des-charges/
`-- frontend/   package.json vite.config.js src/App.jsx
\end{lstlisting}

\section{Glossaire}
\begin{description}[leftmargin=3.2cm,style=nextline]
\item[CI/CD] Intégration continue et livraison/déploiement continus.
\item[Rollback] Retour à la dernière version saine déployée.
\item[Pull request] Proposition de modification de code soumise à relecture.
\item[LLM] Grand modèle de langage.
\item[Dry-run] Mode simulation : l'action est décrite mais pas exécutée.
\item[Abstention] Refus de conclure quand le signal est insuffisant.
\end{description}

\end{document}
